\chapter*{Lời Tựa}
\addcontentsline{toc}{chapter}{Lời Tựa}
\markboth{Lời Tựa}{Lời Tựa}

\begin{truyen}{Lời Tựa}{The Classroom Multiverse}
\chuhoa{L}{ớp học hôm nay} không còn là một đường thẳng từ bục giảng xuống bàn học. Nó giống một ngã tư hơn: học sinh mang theo thuật ngữ chữa lành, phụ huynh mang theo kỳ vọng, thầy cô mang theo trách nhiệm và cả những lúc nóng nảy, còn mạng xã hội thì lặng lẽ đứng phía sau, bơm thêm ngôn ngữ và thái độ vào mọi cuộc đối đáp.
\textit{(Today's classroom is no longer a straight line from podium to desk. It is more like an intersection: students bring healing language, parents bring expectations, teachers bring responsibility and occasional impatience, and social media quietly stands behind everything, feeding new words and attitudes into every exchange.)}

Quyển XXXIII này chọn nhìn lớp học như một \textbf{đa giác học đường}. Mỗi cạnh là một lực kéo: học sinh muốn được hiểu, giáo viên muốn giữ nề nếp, phụ huynh muốn bảo vệ con, còn tương lai thì luôn gọi tên những ảo tưởng rất hấp dẫn. Khi các cạnh ấy chạm nhau, tiếng cười bật ra. Nhưng nếu nhìn sâu hơn, bên dưới tiếng cười là những vấn đề rất thật.
\textit{(This volume chooses to view school as a \textbf{classroom polygon}. Each edge is a pulling force: students want understanding, teachers want order, parents want to protect their children, and the future keeps calling with attractive illusions. When those edges collide, laughter appears. Yet beneath that laughter are very real issues.)}

Cuốn này vẫn giữ cái chất chát, nhanh, gọn của dòng sách học đường trước. Nhưng nó mở rộng góc nhìn hơn: không chỉ châm biếm học sinh, mà còn soi lại người lớn. Vì muốn giáo dục tử tế hơn, đôi khi người lớn cũng phải chịu bị phản biện.
\textit{(This book keeps the sharp, quick, compact tone of the earlier school volumes. But it widens the lens: it does not only satirize students, it also examines adults. If education is to become kinder and wiser, adults must sometimes accept being challenged too.)}
\end{truyen}

\section*{Vì Sao Quyển Này Phải Rộng Hơn?}

\begin{truyen}{Không Còn Chỉ Là Chuyện Của Một Phía}{Beyond One-Sided Blame}
Nhiều quyển sách học đường trước đây có thể đứng hẳn về một phía để tạo tiếng cười nhanh: hoặc học sinh lý sự, hoặc thầy cô nóng nảy, hoặc phụ huynh áp đặt. Cách viết ấy có ích khi ta muốn chỉ ra một thói xấu rõ ràng. Nhưng lớp học bây giờ phức tạp hơn. Một học sinh im lặng có thể vừa lười, vừa kiệt sức, vừa sợ bị chê. Một giáo viên nghiêm có thể vừa đúng nguyên tắc, vừa thiếu kỹ năng nói để người nghe không thấy bị dồn ép. Một phụ huynh bênh con có thể vừa thương con thật, vừa sợ cảm giác mình đã thất bại trong việc nuôi dạy.
\textit{(Earlier school books could stand firmly on one side to create quick satire: students making excuses, teachers losing temper, or parents forcing expectations. That works when one flaw is clear. But today's classroom is more complex. A silent student may be lazy, exhausted, and afraid of being judged at the same time. A strict teacher may be principled yet lack the language to correct without cornering. A parent defending a child may be loving that child sincerely while also fearing their own failure.)}

Vì vậy, Quyển XXXIII không cố dựng thêm một phiên tòa. Nó cố dựng một tấm gương nhiều mặt. Người đọc có thể soi thấy người khác trước, nhưng rồi sẽ phải soi tới mình. Đó là lý do cuốn này cần nhiều lớp hơn, chậm hơn ở vài chỗ, và không ngại đi từ tiếng cười sang những đoạn trầm.
\textit{(Therefore, Volume XXXIII does not try to build another courtroom. It tries to build a many-sided mirror. Readers may first see others in it, but eventually they must see themselves. That is why this book needs more layers, a slower pace in places, and a willingness to move from satire to quieter reflection.)}
\end{truyen}

\section*{Người Viết Muốn Giữ Điều Gì?}

\begin{baihoc}
Nếu chỉ châm biếm mà không mở đường sửa, cuốn sách sẽ làm người đọc hả hê nhưng không trưởng thành hơn. Nếu chỉ vuốt ve mà không dám chỉ thẳng cái dở, cuốn sách sẽ tử tế theo kiểu vô dụng. Quyển này không chọn làm vừa lòng; quyển này chọn nói thật đủ đau để người đọc còn nhớ mà sửa.
\textit{(If a book only satirizes without opening a path to change, it may entertain readers but not help them grow. If it only comforts without pointing directly at what is flawed, it becomes kind in a useless way. This volume does not choose to please; it chooses to tell a truth sharp enough to stay in the reader's mind.)}
\end{baihoc}

\begin{goigiaovien}
\textbf{Gợi ý đọc lời tựa trên lớp:} giáo viên có thể chọn một đoạn ngắn trong phần này để mở đầu tiết sinh hoạt năm học mới, rồi mời học sinh nói xem các em thấy lớp mình hiện nay đang vướng vào ``cạnh'' nào của đa giác nhiều nhất: áp lực thành tích, sĩ diện bạn bè, khoảng cách với cha mẹ, hay ảo tưởng tương lai.
\end{goigiaovien}

\begin{baihoc}
Khi mọi bên đều có một phần đúng và một phần quá lời, điều cần nhất không phải là thắng nhau, mà là nhìn rõ điều gì đang làm nhau mệt.
\textit{(When every side is partly right and partly excessive, the most important thing is not winning, but seeing what is exhausting everyone.)}
\end{baihoc}

\begin{ghinhoanh}
\textbf{Short English to remember:}

The classroom is not one voice.

It is a collision of expectations.
\end{ghinhoanh}

\begin{tuhocnhanh}
1. Trong lớp học, em thường thấy mình đứng ở cạnh nào của ``đa giác'': học sinh, giáo viên, phụ huynh, hay người quan sát? \textit{(Which edge of the polygon do you usually stand on: student, teacher, parent, or observer?)}

2. Em đang có xu hướng đòi người khác hiểu mình, hay cố hiểu người khác? \textit{(Are you asking others to understand you, or trying to understand them?)}

3. Nếu phải chọn một điều để bớt cực đoan hơn trong cách nhìn về trường học, em sẽ bớt trách ai trước? \textit{(If you had to become less extreme in one school-related judgment, whom would you blame less first?)}
\end{tuhocnhanh}

\begin{flushright}
\textbf{Nguyễn Văn Sang}\\
\textit{GV THPT Nguyễn Hữu Cảnh - TP HCM}
\end{flushright}

\ngancach